\chapter*{Summary}
\addcontentsline{toc}{chapter}{Summary}
\markboth{Summary}{Summary}

\begingroup
\small

This thesis investigates whether a small open-weight language model can be
adapted to distinguish unwanted email from legitimate correspondence.
Qwen3-0.6B was fine-tuned with Low-Rank Adaptation (LoRA). The classifier used
only the message subject and body. The positive class included advertising
spam, phishing, and fraudulent messages. Attachments, transport headers,
sender reputation, and email authentication results were outside the scope of
the experiments.

Public email datasets were converted to a common schema, cleaned, deduplicated,
and split into training, validation, and test subsets in an 80/10/10 ratio.
Samples from three additional datasets were used for external validation. Sixteen
configurations were compared for classification based on the probability of
the next token. They differed in context length, learning rate, LoRA rank,
scaling factor, and dropout. Intermediate training checkpoints were evaluated.
Three alternative task formulations were also tested: sequence classification,
label generation with parsing, and structured response generation.

The unmodified model showed a strong preference for the spam class and obtained
final scores of 57.61\% for next-token scoring and 57.30\% for response
generation. The configuration with the highest external-validation score used
a context of 512 tokens, a learning rate of \(10^{-4}\), LoRA rank \(r=32\),
scaling factor \(\alpha=64\), and dropout of 0.05. Its aggregate score was
98.0\%.
Next-token scoring was selected because it derives the decision directly from
the probabilities of the \texttt{spam} and \texttt{ham} labels, whereas the
generative variants require response parsing.

The final comparison used four test sets containing 2,000 messages each. The
structured response method obtained the highest aggregate score, 97.27\%,
while next-token scoring reached 97.24\%. The 95\% confidence interval for the
difference included zero, and the observed difference was small relative to its
uncertainty. Among the baseline methods,
DistilBERT obtained 93.20\%, followed
by TF-IDF with Naive Bayes at 91.82\% and SVM at 90.98\%.

On a common sample of 1,000 messages and an NVIDIA A100 GPU, Qwen3 processed
125.1 messages per second and used approximately 2,595 MiB of RAM and 2,787
MiB of GPU memory. DistilBERT reached 283.5 messages per second, while the
TF-IDF classifiers processed between 1.3 and 6.4 thousand messages per second.

Production deployment would require evaluation on current email traffic,
repeated training runs, and measurements under realistic server load.
Qwen3-0.6B could be evaluated as a second-stage classifier for messages not
resolved by authentication checks, sender reputation, rules, or a less
expensive text classifier.

\medskip
\noindent\textbf{Keywords:} unwanted email, spam detection, phishing, language
models, Qwen3, LoRA, text classification
\endgroup
